% GENERATED FILE. Edit canonical lessons, then run npm run generate:books.
% canonical-source-hash: fnv1a64:5a4dbc36c3429bc6
% canonical-lessons: GU-C245-grhkarya, GU-C245-jodni, GU-C245-bhusvu, GU-C245-rabar, GU-C245-uccarvu

\chapter{Homework and spelling}
\label{ch:gu-homework-and-spelling}
\hlchaptermodality{245}

\begin{chapteropening}
\textbf{By the end of this chapter:} \emph{I can talk about homework, spelling and saying words clearly.}

Complete the last lesson of chapter 245: I can talk about homework, spelling and saying words clearly.
\end{chapteropening}

\section[\texorpdfstring{\gu{ગૃહકાર્ય} (gṛhkārya)}{ગૃહકાર્ય (gṛhkārya)}]{\gu{ગૃહકાર્ય} (gṛhkārya) — homework}

\label{lesson:GU-C245-grhkarya}



\begin{quote}
Before the new one: say the Gujarati for to explain, then the Gujarati for to solve.
\end{quote}

\subsection*{You'll want to know: \gu{ગૃહકાર્ય}}
\textbf{\gu{ગૃહકાર્ય}} — \emph{gṛhkārya} — \textquotedblleft{}homework\textquotedblright{}.

\begin{quote}
\textbf{\gu{હું} \gu{સાંજે} \gu{ગૃહકાર્ય} \gu{કરું} \gu{છું}.} — \emph{hũ sā̃je gṛhkārya karũ chũ} — \textquotedblleft{}I do my homework in the evening.\textquotedblright{}
\end{quote}

A neuter noun.

\begin{cousinweb}
\textbf{\gu{ગૃહકાર્ય}} (\emph{gṛhkārya}) is a modern learned compound built from two tatsamas: Sanskrit \emph{gṛha}, \textquotedblleft{}house\textquotedblright{} (CDIAL 4240), and \emph{kārya}, \textquotedblleft{}what is to be done, business\textquotedblright{} (CDIAL 3078).
\end{cousinweb}

\begin{grammarlens}[title={What you've built}]
One more word for everyday life at A2.
\end{grammarlens}

\subsection*{Your turn}
\begin{itemize}
\raggedright
  \item \emph{Say it:} \emph{gṛhkārya}
  \item \emph{Say it:} \emph{gṛhkārya}, once more
  \item \emph{Say it:} \emph{hũ sā̃je gṛhkārya karũ chũ} — I do my homework in the evening.
  \item \emph{Recall:} say the Gujarati for to explain, then the Gujarati for to solve, then say \emph{gṛhkārya} again
\end{itemize}

\begin{tcolorbox}[breakable,colback=teal!4,colframe=teal!35!black,title={Before you move on}]
What does \gu{ગૃહકાર્ય} mean? (\textquotedblleft{}Homework\textquotedblright{}.) Say it once more.
\end{tcolorbox}
\section[\texorpdfstring{\gu{જોડણી} (joḍṇī)}{જોડણી (joḍṇī)}]{\gu{જોડણી} (joḍṇī) — spelling}

\label{lesson:GU-C245-jodni}



\begin{quote}
Before the new one: say the Gujarati for to solve, then the Gujarati for homework.
\end{quote}

\subsection*{You'll want to know: \gu{જોડણી}}
\textbf{\gu{જોડણી}} — \emph{joḍṇī} — \textquotedblleft{}spelling\textquotedblright{}.

\begin{quote}
\textbf{\gu{આ} \gu{શબ્દની} \gu{જોડણી} \gu{અઘરી} \gu{છે}.} — \emph{ā śabdanī joḍṇī agharī che} — \textquotedblleft{}The spelling of this word is hard.\textquotedblright{}
\end{quote}

\begin{cousinweb}
\textbf{\gu{જોડણી}} (\emph{joḍṇī}), the joining of letters, is built on the inherited verb \emph{joḍvũ}, \textquotedblleft{}to join\textquotedblright{}, with the suffix \emph{-aṇī}. CDIAL lists \emph{joḍvũ} under \emph{*yōṭayati}, \textquotedblleft{}joins\textquotedblright{}, a reconstructed form in its entry \emph{*yuṭati}, \textquotedblleft{}is joined\textquotedblright{}, and calls the history of the group obscure (CDIAL 10496).
\end{cousinweb}

\begin{grammarlens}[title={What you've built}]
One more word for everyday life at A2.
\end{grammarlens}

\subsection*{Your turn}
\begin{itemize}
\raggedright
  \item \emph{Say it:} \emph{joḍṇī}
  \item \emph{Say it:} \emph{joḍṇī}, once more
  \item \emph{Say it:} \emph{ā śabdanī joḍṇī agharī che} — The spelling of this word is hard.
  \item \emph{Recall:} say the Gujarati for to solve, then the Gujarati for homework, then say \emph{joḍṇī} again
\end{itemize}

\begin{tcolorbox}[breakable,colback=teal!4,colframe=teal!35!black,title={Before you move on}]
What does \gu{જોડણી} mean? (\textquotedblleft{}Spelling\textquotedblright{}.) Say it once more.
\end{tcolorbox}
\section[\texorpdfstring{\gu{ભૂંસવું} (bhū̃svũ)}{ભૂંસવું (bhū̃svũ)}]{\gu{ભૂંસવું} (bhū̃svũ) — to rub out, to erase}

\label{lesson:GU-C245-bhusvu}



\begin{quote}
Before the new one: say the Gujarati for homework, then the Gujarati for spelling.
\end{quote}

\subsection*{You'll want to know: \gu{ભૂંસવું}}
\textbf{\gu{ભૂંસવું}} — \emph{bhū̃svũ} — \textquotedblleft{}to rub out, to erase\textquotedblright{}.

\begin{quote}
\textbf{\gu{તે} \gu{ખોટો} \gu{શબ્દ} \gu{ભૂંસે} \gu{છે}.} — \emph{te khoṭo śabda bhū̃se che} — \textquotedblleft{}He rubs out the wrong word.\textquotedblright{}
\end{quote}

\begin{cousinweb}
The origin of \textbf{\gu{ભૂંસવું}} (\emph{bhū̃svũ}) is not settled. CDIAL lists no Gujarati verb of this form and meaning, and our other sources do not trace it.
\end{cousinweb}

\begin{grammarlens}[title={What you've built}]
One more word for everyday life at A2.
\end{grammarlens}

\subsection*{Your turn}
\begin{itemize}
\raggedright
  \item \emph{Say it:} \emph{bhū̃svũ}
  \item \emph{Say it:} \emph{bhū̃svũ}, once more
  \item \emph{Say it:} \emph{te khoṭo śabda bhū̃se che} — He rubs out the wrong word.
  \item \emph{Recall:} say the Gujarati for homework, then the Gujarati for spelling, then say \emph{bhū̃svũ} again
\end{itemize}

\begin{tcolorbox}[breakable,colback=teal!4,colframe=teal!35!black,title={Before you move on}]
What does \gu{ભૂંસવું} mean? (\textquotedblleft{}To rub out, to erase\textquotedblright{}.) Say it once more.
\end{tcolorbox}
\section[\texorpdfstring{\gu{રબર} (rabar)}{રબર (rabar)}]{\gu{રબર} (rabar) — eraser, rubber}

\label{lesson:GU-C245-rabar}



\begin{quote}
Before the new one: say the Gujarati for spelling, then the Gujarati for to rub out.
\end{quote}

\subsection*{You'll want to know: \gu{રબર}}
\textbf{\gu{રબર}} — \emph{rabar} — \textquotedblleft{}eraser, rubber\textquotedblright{}.

\begin{quote}
\textbf{\gu{મારું} \gu{રબર} \gu{ક્યાં} \gu{છે}?} — \emph{mārũ rabar kyā̃ che?} — \textquotedblleft{}Where is my eraser?\textquotedblright{}
\end{quote}

\begin{cousinweb}
\textbf{\gu{રબર}} (\emph{rabar}) is a loan from English \emph{rubber}, the material and the eraser made of it.
\end{cousinweb}

\begin{grammarlens}[title={What you've built}]
One more word for everyday life at A2.
\end{grammarlens}

\subsection*{Your turn}
\begin{itemize}
\raggedright
  \item \emph{Say it:} \emph{rabar}
  \item \emph{Say it:} \emph{rabar}, once more
  \item \emph{Say it:} \emph{mārũ rabar kyā̃ che?} — Where is my eraser?
  \item \emph{Recall:} say the Gujarati for spelling, then the Gujarati for to rub out, then say \emph{rabar} again
\end{itemize}

\begin{tcolorbox}[breakable,colback=teal!4,colframe=teal!35!black,title={Before you move on}]
What does \gu{રબર} mean? (\textquotedblleft{}Eraser, rubber\textquotedblright{}.) Say it once more.
\end{tcolorbox}
\section[\texorpdfstring{\gu{ઉચ્ચારવું} (uccārvũ)}{ઉચ્ચારવું (uccārvũ)}]{\gu{ઉચ્ચારવું} (uccārvũ) — to pronounce, to utter}

\label{lesson:GU-C245-uccarvu}



\begin{quote}
Before the new one: say the Gujarati for to rub out, then the Gujarati for eraser.
\end{quote}

\subsection*{You'll want to know: \gu{ઉચ્ચારવું}}
\textbf{\gu{ઉચ્ચારવું}} — \emph{uccārvũ} — \textquotedblleft{}to pronounce, to utter\textquotedblright{}.

\begin{quote}
\textbf{\gu{તે} \gu{દરેક} \gu{શબ્દ} \gu{ધીમે} \gu{ઉચ્ચારે} \gu{છે}.} — \emph{te darek śabda dhīme uccāre che} — \textquotedblleft{}She pronounces every word slowly.\textquotedblright{}
\end{quote}

\begin{cousinweb}
\textbf{\gu{ઉચ્ચારવું}} (\emph{uccārvũ}) comes from Sanskrit \emph{uccārayati}, \textquotedblleft{}utters\textquotedblright{}, the causative of \emph{úccarati}. CDIAL lists the inherited Gujarati \emph{ucārvũ}, \textquotedblleft{}to pronounce, speak\textquotedblright{}; the written \emph{cc} follows the Sanskrit, a semi-tatsama (CDIAL 1641).
\end{cousinweb}

\begin{grammarlens}[title={What you've built}]
One more word for everyday life at A2.
\end{grammarlens}

\subsection*{Your turn}
\begin{itemize}
\raggedright
  \item \emph{Say it:} \emph{uccārvũ}
  \item \emph{Say it:} \emph{uccārvũ}, once more
  \item \emph{Say it:} \emph{te darek śabda dhīme uccāre che} — She pronounces every word slowly.
  \item \emph{Recall:} say the Gujarati for to rub out, then the Gujarati for eraser, then say \emph{uccārvũ} again
\end{itemize}

\begin{tcolorbox}[breakable,colback=teal!4,colframe=teal!35!black,title={Before you move on}]
What does \gu{ઉચ્ચારવું} mean? (\textquotedblleft{}To pronounce, to utter\textquotedblright{}.) Say it once more.
\end{tcolorbox}
